The previous section establishes that properly initialised cohorts sit in
the floor-zero regime, where all observed worst-task excess is in
principle recoverable. That is an invitation to build an encoder that
recovers it, and we did: a rate-distortion encoder using incoherent
orthogonal mixing and scalar quantization, regularised with a single
Tikhonov ridge term (\S\ref{sec:achv}, implementation in
Appendix~\ref{app:manifest}). On the shared-initialisation cohorts our
earlier draft measured, it wins comfortably against every baseline.

This section reports what happened when we tested that result properly.
Every decision rule quoted below was committed to version control before
the corresponding compute ran, and the commit ordering is the audit trail.
The pre-registrations are reproduced verbatim in
Appendix~\ref{app:prereg}.

\subsection{The encoder's advantage does not survive}
\label{subsec:q1}

\paragraph{Pre-registered question.} Does the encoder's win over the
best baseline survive the move from the shared-initialisation cohort to
independently initialised ones? The rule, fixed in advance: for each base
model, compare the encoder against the champion baseline selected on the
mean over the three independent cohorts, with a tie threshold of $0.005$
nats and a one-directional noise gate at $2\times\mathrm{SE}$ that may
only downgrade a win or loss to a tie, never promote.

\begin{table}[t]
\centering
\small
\caption{Pre-registered replication on independently initialised cohorts.
Worst-task NLL excess in nats, mean over three cohorts. Positive $d$
favours the encoder. Per-cell standard deviation across cohorts is
$0.0001$ to $0.0034$, so the noise gate downgraded nothing.}
\label{tab:q1}
\begin{tabular}{lccccc}
\toprule
base & rd-encoder ridge & champion baseline & $d$ & $2\times\mathrm{SE}$ & result \\
\midrule
Llama-3.1-8B & 0.0717 & 0.1052 \; (TVQ$_2$) & $+0.0335$ & 0.0010 & wins \\
Mistral-7B   & 0.0597 & 0.0543 \; (TIES)    & $-0.0055$ & 0.0024 & loses \\
Qwen2.5-7B   & 0.0143 & 0.0141 \; (TIES)    & $-0.0002$ & 0.0004 & ties \\
Yi-1.5-9B    & 0.0637 & 0.0487 \; (TIES)    & $-0.0150$ & 0.0011 & loses \\
\bottomrule
\end{tabular}
\end{table}

\paragraph{Outcome: it does not survive.} Table~\ref{tab:q1} gives one
win, one tie and two losses. Both pre-specified robustness lines agree.
Re-selecting the champion inside each cohort, the variant deliberately
biased against our method, returns the same one/one/two. Applying the
original single-cohort rule three times returns a weakened verdict on
\texttt{indep1} but a failure on \texttt{indep2} and \texttt{indep3}, so
the earlier optimistic reading was the best of three draws.

The blunt summary is that \textbf{TIES, a 2023 baseline, is the best
method on three of four base models} once the adapters are initialised
properly. Our encoder is beaten on Mistral and Yi, ties on Qwen, and wins
only on Llama. The Llama win is not a rank artifact: rank-matching the
encoder to the baselines still leaves it at $0.0795$ against the
champion's $0.1052$.

We think the honest reading is that the encoder's apparent advantage was
substantially a property of the degenerate cohort it was developed on.
In the floor-positive regime there is a large irreducible component that
different methods approach differently, and a method tuned on that regime
need not generalise to the regime where the floor is gone.

\subsection{A claim we withdraw}
\label{subsec:withdrawn}

The result in \S\ref{subsec:q1} raises an obvious and much more
interesting possibility: if method rankings shift when initialisation
changes, then published merging benchmarks, which as far as we can tell
mostly use shared-seed cohorts, might be measuring initialisation
geometry rather than method quality. We described this on first sight as
the strongest result available to us.

We pre-registered a control to test it, and the control did not support
the claim. The control asks whether the method ranking is stable
\emph{within} the independent regime, across the three cohorts that
differ only in initialisation draw. If rankings are unstable within a
single regime, then ranking changes across regimes cannot be attributed
to the regime. Two of four base models had a non-unanimous top-1 method
across the three cohorts, above the threshold fixed in advance, so by our
own rule the attribution fails and \textbf{the claim that published
merging benchmarks are confounded by shared-initialisation geometry is
withdrawn}.

We record one caveat, and we deliberately did not let it reverse the
verdict. The control as written counts a top-1 change as instability
regardless of margin, and both flips involve methods separated by far less
than the tie threshold: on Qwen the top two are $0.0140$ and $0.0141$, a
coin flip. A margin-aware version of the control would very likely give a
different answer. But that version was not what we pre-registered, and
rewriting a decision rule after seeing which way it fell is exactly the
practice the pre-registration exists to prevent. We state the defect,
leave the verdict, and flag a margin-aware control as future work.

What survives is weaker and still worth saying: single-seed method
rankings on this metric are not reproducible, so merging comparisons
should report multiple initialisation draws.

\subsection{The predicted rate exponent is falsified on real adapters}
\label{subsec:w3}

Theorem~\ref{thm:general} predicts distortion decaying as
$2^{-2R/\deff}$, which on a log-rate axis is a slope in a narrow band. We
pre-registered the band as $[-2.4, -1.6]$, to hold on at least three of
four base models, and swept quantizer width $b \in \{1,2,3,4,8,16\}$.

It fails. Of the two base models where the curve can be fitted at all,
the measured slopes are $-0.10$ and $-0.21$, roughly an order of
magnitude off the predicted band, so the criterion is met on zero of the
fittable models. On the other two the fit is undefined because the excess
at finite rate falls \emph{below} its own $b = \infty$ value, which the
model does not allow.

The data instead show something flat and operationally useful: with the
ridge on, everything from $b = 2$ upward is indistinguishable. Two bits
buys the entire effect and sixteen bits adds nothing. We report this as an
empirical regularity, not as a confirmation of the theory, because it is
not one. The theory's floor term is what our measurements support; its
rate term is not, at least at the ranks and scales reachable here.

A caveat on the falsification itself: this sweep was run on a single
initialisation draw. A single draw is thin evidence for a falsification,
though the flatness is far too large and too consistent across four base
models to be plausibly explained by initialisation noise.

\subsection{DARE: a negative result with a mechanism, and its limits}
\label{subsec:dare}

DARE~\citep{yu2024dare} drops a random fraction of task-vector entries
and rescales the survivors by $1/(1-p)$. Composed with task arithmetic it
does essentially nothing in our setting: across a density sweep from
$0.05$ to $0.5$ on four base models, it beats plain task arithmetic in
three of twenty cells and all three are within $0.0005$ nats, a fifth of
our tie threshold.

Unlike most negative results this one has a mechanism that predicts its
own shape. The rescaling makes the masked task vector an unbiased
estimator of the unmasked one. For a loss convex in the merged delta,
Jensen's inequality then gives
$\E[L(D_{\mathrm{DARE}})] \geq L(D_{\mathrm{TA}})$, so the composition
\emph{cannot} beat task arithmetic in expectation, and to second order the
penalty is $\tfrac12\tr(H\Sigma)$ with $\Sigma \propto (1-p)/p$, which
predicts monotone degradation as density falls. That is what we observe,
on three of four base models. We verified our implementation against the
reference one and confirmed the operator reproduces the analytic signature
$\sqrt{(1-p)/p}$ to three decimals, so this is a property of the method
and not of our code.

\paragraph{The negative result does not generalise, and we tested that.}
The Jensen argument applies only to the task-arithmetic composition, since
TIES is biased and the mask changes which coordinates survive trimming. We
pre-registered the DARE-TIES composition test before implementing it. The
outcome is genuinely mixed: DARE composed with TIES \emph{helps}
decisively on Llama ($+0.0462$ nats, against a gate of $0.0036$) and hurts
decisively on the other three. By the rule fixed in advance this is a
mixed verdict and neither direction is claimed.

So the correct scope for the negative result is DARE composed with task
arithmetic, not DARE. We flag that Llama is now the odd model out in two
separate pre-registered tests, here and in Table~\ref{tab:q1}, without
proposing a mechanism for it; a base-level moderator that neither
pre-registration anticipated deserves a designed experiment rather than
another post-hoc reading.
